% =============================================================================
% INTRODUCTION
% Status: DRAFT - Framework established, needs literature review
% =============================================================================

\subsection{Background}

Unmanned Aerial Vehicles (UAVs) have become essential tools in precision agriculture, enabling high-resolution monitoring of crop health, irrigation patterns, and field conditions. Modern drone surveys generate large volumes of imagery---often hundreds to thousands of images per flight session---creating significant data management challenges.

\todo{Add citations: drone imagery in agriculture, data volume challenges}

A typical agricultural drone survey may include:
\begin{itemize}
    \item \textbf{Systematic mapping flights}: High-altitude images for orthomosaic generation
    \item \textbf{Targeted inspection flights}: Low-altitude close-ups of specific crop areas
    \item \textbf{Oblique/transitional images}: Various angles during ascent/descent
\end{itemize}

Organizing this heterogeneous imagery typically requires manual review or relies solely on GPS metadata, which captures \emph{where} images were taken but not \emph{what} they contain semantically.

\subsection{Vision-Language Models for Image Understanding}

Recent advances in vision-language models, particularly CLIP (Contrastive Language-Image Pre-training), have demonstrated remarkable zero-shot capabilities for image understanding. These models learn joint embeddings of images and text, enabling semantic similarity comparisons without task-specific training.

\observation{CLIP embeddings capture semantic content (altitude/purpose) rather than just spatial position---this is a key distinction from GPS-based organization}

\subsection{Research Gap}

While vision-language models have been extensively studied for general image classification and retrieval, their application to agricultural drone imagery organization remains underexplored. Specifically:

\begin{enumerate}
    \item Can pretrained models (without agricultural fine-tuning) meaningfully cluster drone imagery?
    \item What semantic distinctions do these models capture in agricultural contexts?
    \item How do embedding-based clusters compare to traditional metadata-based organization?
\end{enumerate}

\todo{Expand literature review: existing work on drone image organization, CLIP applications in remote sensing}
